% ============================================================
%  Needs Expression / Problem Statement
%  Project: AI-Based IDS/IPS — Port Scan Detection
%  Module: Artificial Intelligence — Semester 2
% ============================================================
\documentclass[12pt,a4paper]{article}

% --- Packages ---
\usepackage[utf8]{inputenc}
\usepackage[T1]{fontenc}
\usepackage[english]{babel}
\usepackage{geometry}
\usepackage{enumitem}
\usepackage{booktabs}
\usepackage{array}
\usepackage{xcolor}
\usepackage{titlesec}
\usepackage{fancyhdr}
\usepackage{hyperref}

\geometry{margin=2.5cm}

% --- Colors ---
\definecolor{sectionblue}{RGB}{15, 23, 42}
\definecolor{accentcyan}{RGB}{6, 182, 212}

% --- Section styling ---
\titleformat{\section}
  {\large\bfseries\color{sectionblue}}
  {\thesection.}{0.5em}{}[\vspace{-0.5em}\textcolor{accentcyan}{\rule{\textwidth}{0.4pt}}]

\titleformat{\subsection}
  {\normalsize\bfseries\color{sectionblue}}
  {\thesubsection.}{0.5em}{}

% --- Header/Footer ---
\pagestyle{fancy}
\fancyhf{}
\fancyhead[L]{\small\textcolor{gray}{Needs Expression --- IDS/IPS Port Scan}}
\fancyhead[R]{\small\textcolor{gray}{AI Module --- S2}}
\fancyfoot[C]{\small\thepage}
\renewcommand{\headrulewidth}{0.4pt}

% --- Table column types ---
\newcolumntype{L}[1]{>{\raggedright\arraybackslash}p{#1}}
\newcolumntype{C}[1]{>{\centering\arraybackslash}p{#1}}

\begin{document}

% ============================================================
%  TITLE BLOCK
% ============================================================
\begin{center}
  {\Large\bfseries Needs Expression / Problem Statement}\\[0.6em]
  {\large AI-Based Intrusion Detection System}\\[0.3em]
  {\large Detection of \textit{Port Scan \& Network Reconnaissance} Attacks}\\[1.2em]
  \begin{tabular}{rl}
    \textbf{Module:}  & Artificial Intelligence --- Semester 2 \\
    \textbf{Team:}    & 5 members (Engineering --- ENSA) \\
    \textbf{Date:}    & April 2026 \\
  \end{tabular}
\end{center}

\vspace{1em}

% ============================================================
%  PARAGRAPH 1 --- Context / Current Situation
% ============================================================
\section{Context and Current Situation}

In the field of network security, the early detection of reconnaissance
activities represents a critical challenge. \textit{Port Scanning} --- the
technique by which an attacker systematically probes the ports of a network to
identify active hosts, exposed services, and exploitable vulnerabilities ---
constitutes the first phase of virtually all cyberattacks documented under the
\textit{Cyber Kill Chain} framework.

Currently, network infrastructures rely primarily on static rule-based firewalls
and signature-based Intrusion Detection Systems (IDS) such as Snort and
Suricata. These tools operate through known pattern matching and apply fixed
thresholds (e.g., ``more than 100 SYN connections within 10 seconds'').

This approach gives rise to several major shortcomings:

\begin{itemize}[nosep,leftmargin=1.5em]
  \item \textbf{Blindness to slow and stealthy scans:} attackers deliberately
        space their probes (1 packet/second or less) to remain below detection
        thresholds, rendering static rules ineffective.
  \item \textbf{High false positive rate:} legitimate traffic (software updates,
        service discovery, monitoring scans) frequently triggers the same
        signatures as malicious scans, generating a volume of irrelevant alerts
        that exhausts network administrators.
  \item \textbf{Lack of behavioral profiling:} current systems treat each packet
        in isolation, without analyzing the temporal and statistical patterns of
        traffic per source IP over sliding time windows.
  \item \textbf{Inability to classify scan type:} traditional IDS cannot
        distinguish a \textit{SYN Stealth Scan} from a \textit{UDP Scan}, an
        \textit{ACK Scan}, or an \textit{XMAS Scan}, depriving the administrator
        of essential tactical information.
\end{itemize}

These limitations directly impact Security Operations Center (SOC) teams in
terms of response time, alert reliability, and the ability to contain an attack
before it progresses to more destructive phases (exploitation, lateral movement,
exfiltration).

% ============================================================
%  PARAGRAPH 2 --- Problem Statement / Needs Expression
% ============================================================
\section{Problem Statement and Needs Expression}

Given these observations, it becomes necessary to rethink the detection
mechanism in order to \textbf{overcome the limitations of static signature-based
approaches} and equip the network with an intelligent, adaptive, real-time
analytical capability.

The primary objective is to design a system capable of automatically learning the
behavioral patterns of network traffic --- both normal and malicious --- in order
to autonomously detect reconnaissance activities, including slow and stealthy
variants.

In this context, several questions arise:

\begin{enumerate}[nosep,leftmargin=1.5em]
  \item Which network traffic \textit{features} effectively discriminate a port
        scan from legitimate behavior?
  \item How can a machine learning model detect scans distributed over time,
        where static thresholds fail?
  \item What labeled and representative data must be provided to the model to
        ensure reliable classification with a minimal false positive rate?
\end{enumerate}

\vspace{0.5em}
\noindent The problem statement can thus be formulated as follows:

\begin{quote}
\itshape
``How can we design, train, and evaluate a Machine Learning model capable of
classifying network traffic with high reliability to detect
\textbf{Port Scan \& Reconnaissance} attacks, while drastically minimizing the
false positive rate within a controlled network infrastructure?''
\end{quote}

This problem statement aims to deliver a solution that enables: reduction of
detection time, improvement of alert accuracy, and decision support for network
administrators.

% ============================================================
%  PARAGRAPH 3 --- Proposed Solution
% ============================================================
\section{Proposed Solution}

To address this problem, an approach based on \textbf{artificial intelligence},
and more specifically on \textbf{supervised classification learning}, can be
adopted.

The system will exploit labeled network traffic data sourced from recognized
IDS benchmark datasets. By leveraging features such as the number of distinct
ports probed per IP, the ratio of SYN packets without completed handshakes,
flow duration, and the entropy of the targeted port range, the model will learn
to distinguish benign flows from reconnaissance flows.

\vspace{0.5em}
\begin{tabular}{@{}l@{\hspace{1em}}p{11cm}@{}}
  \textbf{Inputs:}      & Feature vectors extracted from network traffic per
                           time window (10\,s and 60\,s) and per source IP:
                           TCP flag counters, distinct port count,
                           inter-arrival time, port entropy, flow duration. \\[0.5em]
  \textbf{Outputs:}      & Binary classification (\textsc{Benign} /
                           \textsc{Port\_Scan}) with confidence score; extended
                           to multi-class scan type classification (SYN / UDP /
                           ACK / XMAS / Connect). \\[0.5em]
  \textbf{Model type:}   & Supervised classification (ensemble decision trees)
                           complemented by an unsupervised anomaly detection
                           layer for slow scans.
\end{tabular}

% ============================================================
%  TWO POSSIBLE AI SOLUTIONS
% ============================================================
\section{Two Proposed AI Solutions}

\subsection{Automation-Oriented Solution}

A \textbf{fully autonomous} system where the classification model, once
deployed, analyzes traffic in real time, issues alerts, and
\textbf{automatically applies} countermeasures (blocking the source IP via
\texttt{iptables}, redirecting to a \textit{honeypot}). No human intervention
is required between detection and response.

\begin{itemize}[nosep,leftmargin=1.5em]
  \item \textit{Advantage:} minimal response time ($<$\,15\,s from detection
        to blocking).
  \item \textit{Risk:} a false positive results in blocking a legitimate user
        without human validation.
\end{itemize}

\subsection{Decision-Support-Oriented Solution}

A system that analyzes traffic and \textbf{presents results on an interactive
dashboard} (port heat map, attacker profile, alert feed with confidence score
and severity level). The administrator \textbf{manually decides} whether to
block the suspicious IP.

\begin{itemize}[nosep,leftmargin=1.5em]
  \item \textit{Advantage:} full human control; risk of blocking due to false
        positives is eliminated.
  \item \textit{Risk:} response time depends on administrator availability.
\end{itemize}

% ============================================================
%  DATA DICTIONARY
% ============================================================
\section{Data Dictionary}

\subsection{Data Sources}

\begin{table}[h!]
\centering
\small
\begin{tabular}{L{3.2cm} L{3.5cm} C{2.8cm} L{3.5cm}}
  \toprule
  \textbf{Dataset} & \textbf{Subset} & \textbf{Approx.\ Volume} & \textbf{Source} \\
  \midrule
  CICIDS2017       & PortScan category       & $\sim$158,000 flows   & UNB / CIC (Canada) \\
  UNSW-NB15        & Reconnaissance category & $\sim$13,000 records  & UNSW (Australia) \\
  \bottomrule
\end{tabular}
\end{table}

\subsection{Input Variables (\textit{Features})}

\begin{table}[h!]
\centering
\small
\begin{tabular}{L{4cm} C{2.5cm} C{2.5cm} L{4cm}}
  \toprule
  \textbf{Variable} & \textbf{Type} & \textbf{Relevance} & \textbf{Role} \\
  \midrule
  Distinct Dst Ports / IP & Quantitative & Critical   & Core scanning signature \\
  SYN Flag Count          & Quantitative & Critical   & SYN stealth scan detection \\
  RST Flag Count          & Quantitative & Very High  & Probe rejection indicator \\
  Flow Duration           & Quantitative & High       & Very short in fast scans \\
  Total Fwd Packets       & Quantitative & High       & Typically 1 per probe \\
  IAT Mean                & Quantitative & High       & Large in slow scans \\
  Port Range Entropy      & Quantitative & High       & Randomized vs sequential scan \\
  ACK Flag Count          & Quantitative & Medium     & ACK scan detection \\
  Unique Dst IPs / Src    & Quantitative & Medium     & Host discovery phase \\
  Bwd Packet Length       & Quantitative & Medium     & RST vs SYN-ACK ratio \\
  TTL Value               & Quantitative & Low        & OS fingerprinting attempts \\
  \bottomrule
\end{tabular}
\end{table}

\subsection{Output Variable (\textit{Label})}

\begin{table}[h!]
\centering
\small
\begin{tabular}{L{4cm} C{2.5cm} L{7cm}}
  \toprule
  \textbf{Variable} & \textbf{Type} & \textbf{Values} \\
  \midrule
  Traffic class (v1) & Qualitative & \textsc{Benign} / \textsc{Port\_Scan} \\
  Scan type (v2)     & Qualitative & SYN / UDP / ACK / XMAS / Connect \\
  \bottomrule
\end{tabular}
\end{table}

\end{document}
